% =====================================================================
% Front-matter rewrite for main_25 (JFM submission)
% British spelling; no em-dashes; numbers to be bound to macros from
% macros_v3.tex / numbers.json before merging (search REVIEW-NUMBER).
% =====================================================================

% ----------------------------- TITLE --------------------------------
% Recommended (en-dash in the compound noun follows Fukami et al. 2024):
\title{Predictive reduced-order states for wall-pressure estimation of
extreme vortex gust--airfoil interactions}
% Alternatives, in order of preference:
% \title{Wall-pressure estimation of extreme vortex gust--airfoil
%   encounters through predictive reduced-order states}
% \title{A wall-estimable predictive reduced state for extreme
%   vortex gust--airfoil interactions}
% Rationale: the physical task (wall-pressure estimation) leads; the
% method class (predictive reduced-order states) follows; "lift-anchored"
% and "JEPA" are results and lineage, not title material.

% Keywords: verify against the current JFM list before submission.
% \keywords{vortex shedding, low-dimensional models, machine learning}

% ---------------------------- ABSTRACT -------------------------------
% ~245 words. Three numbers only. Physical units alongside R^2.
% Restores the attribution sentence (supervision vs predictive objective)
% and the oracle clause; removes all divergence-count claims.
\begin{abstract}
Extreme vortex gusts drive a post-stall airfoil through a transient
reorganisation of its wake, and the release parameters alone do not
determine the flow once the vortex has passed. We ask whether a reduced
state of such an encounter can retain the wake dynamics, remain usable
under a low-order forecast model, and be estimated from the sparse wall
pressure a wing can carry. Direct numerical simulations of a NACA 0012
airfoil at $\alpha = 14^\circ$ and $Re = 5000$, perturbed by Taylor
vortices spanning gust ratio, core diameter and offset, are used to
compare predictive, reconstruction-trained and
proper-orthogonal-decomposition coefficient states at matched dimension.
A controlled supervision study shows that these requirements are
supplied by different ingredients. Under the predictive objective
considered, a scalar lift target is necessary to prevent collapse of the
latent representation, a near-body force-element target sharpens the
peak loads, and a wake target is required to retain the wake enstrophy
and circulation; the multi-step predictive objective contributes
stability under rollout rather than instantaneous readability. Load
information remains accessible at four coefficients, whereas wake-field
reconstruction requires at least sixteen. Under a single direct
multi-horizon forecaster trained identically on every state, the
predictive coordinates are the most forecastable, autoregressive rollout
degrades through error accumulation, and supplying the true gust
parameters provides no benefit. A sequential ensemble filter sensing
eight wall-pressure taps then tracks the lift through the encounter,
with impact-phase $R^2 = 0.794$ (root-mean-square error $0.286$ in lift
units) on the in-distribution test set and $R^2 = 0.837$ at the held-out
$|G| = 4$ boundary; % REVIEW-NUMBER: bind 0.794 / 0.286 / 0.837 to macros
the wake enstrophy is not recovered by the wall-sensor configurations
considered, a limit set by the sensing geometry rather than by the
state.
\end{abstract}

% ------------- FINAL PARAGRAPH OF THE INTRODUCTION -------------------
% Replaces the current four-contribution paragraph (lines 93-120).
% Includes the (i)/(ii)/(iii) requirement list and the roadmap sentence,
% both house style in the lineage corpus. ~185 words; trim the roadmap
% sentence if the editor requires ~120.
Here we ask whether a reduced representation of an extreme vortex-gust
encounter can satisfy three requirements simultaneously: (i) it must
carry the wake variables that distinguish post-impact states beyond the
release parameters; (ii) it must remain usable when advanced by a
low-order forecast model; and (iii) it must be estimable from the sparse
wall pressure a vehicle carries. We compare predictive,
reconstruction-trained and proper-orthogonal-decomposition coefficient
states at matched dimension, on direct numerical simulations of a
post-stall NACA 0012 airfoil perturbed by Taylor vortex gusts.
Controlled ablations separate the roles of the lift, near-body and wake
supervision targets; one direct forecaster, trained identically on every
state, assesses forecastability; and a sequential estimator that senses
only wall pressure tests observability. The requirements prove to be
supplied by different ingredients: observable supervision organises the
latent geometry and its readability, multi-step predictive training
stabilises its forward use, and pressure histories recover the
lift-carrying coordinates but not the complete wake state. The remainder
of the paper is organised as follows. Section 2 describes the flow
configuration, data and endpoints; \S 3 the reduced states, forecaster
and estimators; \S 4 the results, in the order construction (\S 4.1),
compression and forecastability (\S 4.2), wall observability (\S 4.3)
and sequential estimation (\S 4.4); \S 5 discusses the mechanisms and
their limits, and \S 6 concludes.

% --------------------- CONCLUDING REMARKS ----------------------------
% Section title "Concluding remarks" per lineage convention.
% Four short paragraphs: construction, compression, forecasting,
% estimation; then one bounding paragraph. All courtroom language gone.
\section{Concluding remarks}

We asked whether a reduced state of an extreme vortex-gust encounter can
be constructed to carry the wake dynamics, advanced in time, and
recovered from wall pressure alone. On construction: under the
predictive objective considered, a scalar lift target is necessary for a
non-collapsed latent representation, a near-body force-element target
sharpens the peak loads on top of it (peak-region $R^2 = 0.86$ over
three seeds), % REVIEW-NUMBER
and a wake target is required for the state to carry the wake
observables. The collapse result concerns latent health and the
peak-load result concerns a paired increment over an already
lift-supervised baseline; the two rest on different comparisons.

On compression, the statement is two-tier: load information is
approximately unchanged over the tested dimensions under a nonlinear
probe and is linearly accessible at four coefficients, whereas
wake-field reconstruction requires at least sixteen. On forecasting, a
direct multi-horizon forecaster remains usable over forty steps where
autoregressive rollout degrades through error accumulation, and
supplying the true gust parameters provides no benefit, so the
deployable route to the encounter is a model corrected by pressure
rather than the release parameters.

On estimation, a two-stage ensemble filter sensing eight wall-pressure
taps tracks the lift through the encounter, phase-resolved and in
physical units: impact-phase $R^2 = 0.794$ at a root-mean-square error
of $0.286$ lift units on the in-distribution test set, and $R^2 = 0.837$
at the $|G| = 4$ boundary. % REVIEW-NUMBER
Across latent dimension and training seed the predictive states remain
uniformly estimable from the wall, whereas the reconstructive geometry
does not, for a verified observation-side reason; this robustness,
rather than any single configuration, is the deployment argument.

Three results bound the claims. A static delay-embedded inverse matches
the filter's in-distribution impact median, and the filter's advantage
lies in the tails, at the boundary set and in continuous state tracking.
The wake enstrophy is not recovered by the tap counts, delay windows and
estimators considered here, a limit we attribute to the sensing geometry
rather than to the state. And the mid-plane representation shows an
indication of a three-dimensional observability limitation at the
strongest gusts, where the filter's uncertainty consistency degrades
even as the median tracking holds. The natural next steps are a
calibrated process and observation noise model for the filter, and the
same matched-estimator protocol on other parametric flows, to test how
general the estimability ordering reported here is.
